\chapter*{ABSTRAK}
\addcontentsline{toc}{chapter}{ABSTRAK}

\begin{center}

    \textbf{Sistem Informasi Berbasis Machine Learning dan Explainable Artificial Intelligence untuk Estimasi Risiko Hipertensi} \\
    \vspace{0.4cm}
    oleh \\
    \vspace{0.4cm}
    Samuel Franciscus Togar Hasurungan\\
    18222131\\
\end{center}


\begin{spacing}{1.0}
Hipertensi adalah salah satu faktor risiko utama yang berkontribusi terhadap terjadinya penyakit kardiovaskular dengan
prevalensi yang terus meningkat, namun penanganannya masih cenderung reaktif sehingga
banyak kasus baru terdeteksi setelah timbul komplikasi. Pendekatan \textit{machine
learning} berpotensi mengestimasi risiko hipertensi lebih dini, tetapi sifat \textit{black
box} model menghambat kepercayaan klinis, sementara luaran \textit{Explainable Artificial
Intelligence} (XAI) yang berupa nilai numerik sulit dipahami masyarakat awam dan penggunaan
\textit{Large Language Model} (LLM) generik rentan terhadap halusinasi medis. Penelitian ini
bertujuan merancang dan mengimplementasikan sistem estimasi risiko hipertensi yang
mengintegrasikan \textit{machine learning}, XAI, dan LLM agar prediksi tidak hanya akurat,
tetapi juga transparan dan komunikatif. Metode yang digunakan mengikuti kerangka CRISP-DM
dengan data sekunder \textit{Indonesian Family Life Survey} gelombang kelima (IFLS-5), yang
menghasilkan 30.251 responden dewasa dengan delapan fitur non-invasif. Empat algoritma, yaitu
\textit{Logistic Regression}, \textit{Random Forest}, \textit{XGBoost}, dan \textit{CatBoost},
dibandingkan, dan XGBoost dipilih sebagai model akhir dengan ROC-AUC 0,78 serta \textit{recall}
0,71 pada titik operasi \textit{Youden Index}, sehingga memenuhi kriteria keberhasilan berupa
ROC-AUC minimal 0,75 dan \textit{recall} minimal 0,70. Interpretasi model dilakukan menggunakan
SHAP (\textit{SHapley Additive exPlanations}), sedangkan narasi penjelasan dibangun melalui
arsitektur dua lapisan: seluruh angka, arah pengaruh, dan fakta ditetapkan secara deterministik
dari nilai SHAP, sedangkan LLM (\textit{llama-3.1-8b-instant}) hanya bertugas menghaluskan bahasa
tanpa mekanisme temu-kembali eksternal (RAG), sehingga risiko halusinasi dapat ditekan. Seluruh
komponen diintegrasikan ke dalam purwarupa berbasis web. Hasil validasi oleh seorang dokter umum
atas tujuh kasus menunjukkan komponen SHAP memperoleh rerata 3,60 dan komponen narasi LLM 3,93
pada skala Likert 1--5, yang keduanya tergolong valid. Dengan demikian, sistem terbukti mampu
menyajikan estimasi risiko hipertensi yang akurat, transparan, dan mudah dipahami, serta layak
digunakan sebagai alat bantu skrining awal dan edukasi, bukan sebagai pengganti diagnosis tenaga
kesehatan.

Kata kunci: Hipertensi, Machine Learning, Explainable AI, SHAP, Large Language Model.
\end{spacing}
